%-------------------------------------------------------------------------------
% resume.tex -- canonical 1-page resume. This file is PUBLIC: contact details
% that should stay private go in resumes/contact.tex (see `make new-resume`).
% Uses awesome-cv-resume.cls.
% Compile: make resume
%-------------------------------------------------------------------------------
\documentclass[11pt, letterpaper]{awesome-cv-resume}

%-------------------------------------------------------------------------------
% PAGE GEOMETRY
%-------------------------------------------------------------------------------
\geometry{left=1.4cm, top=.8cm, right=1.4cm, bottom=1.0cm, footskip=.5cm}

%-------------------------------------------------------------------------------
% FONTS & COLOURS  (mirror cv.tex settings)
%-------------------------------------------------------------------------------
\fontdir[fonts/]
\colorlet{awesome}{awesome-darknight}
\setbool{acvSectionColorHighlight}{true}

% Make first name bold (overrides the light-weight default in awesome-cv.cls)
\renewcommand*{\headerfirstnamestyle}[1]{{\fontsize{32pt}{1em}\headerfont\bfseries\color{text} #1}}

%-------------------------------------------------------------------------------
% PERSONAL INFORMATION
%-------------------------------------------------------------------------------
\name{Shabnam}{Hakimi, Ph.D.}
% \position{Senior Research Scientist{\enskip\cdotp\enskip}Human-Centered AI}
% \address{San Francisco, CA}
% PRIVATE-CONTACT
\email{mail@shabnamhakimi.com}
\homepage{shabnamhakimi.com}

%-------------------------------------------------------------------------------
\begin{document}

\makecvheader[C]   % centred header, same as cv.tex

%-------------------------------------------------------------------------------
% EXPERIENCE
%-------------------------------------------------------------------------------
\vspace{2mm}
\cvsection{Experience}

\begin{cventries}

  \cventry
    {Senior Research Scientist, Human-Centered AI}
    {Toyota Research Institute}
    {Los Altos, CA}
    {Jan 2021 -- Present}
    {
      \begin{cvitems}
        \item Own end-to-end evaluation frameworks and quality metrics for AI systems (LLMs, VLMs, diffusion models) that combine structured human evaluation (rubrics, benchmark datasets, inter-rater reliability testing) with automated, model-based approaches (LLM-as-judge) to assess model alignment across diverse populations and conditions. Develop new instruments for measuring and modeling unobservable human and artificial agent states for complex phenomena like creativity and motivation.
        \item Lead original research on the causal impact of generative AI interventions on human creativity and decision making. Architect human-aligned AI systems grounded in psychology and neuroscience with rigorous evaluation of downstream behavioral effects.
        \item Build novel multimodal AI and ML methods for consumer demand forecasting. Leverage GAN, VAE, and transformer-based NLP methods for analyzing sparse consumer data. Model behavioral, physiological, and text response data jointly to develop new preference measurement and prediction approaches at scale. Demonstrate improved forecasting accuracy for innovative products under high uncertainty, materially impacting business decision making.
        \item Lead large, multi-site projects with cross-functional teams comprising 5--15 people spanning research, engineering, and product. Define project scope, deliverables, and success metrics; proactively identify risks and develop mitigation strategies. Manage cross-organizational alignment.
        \item Communicate research results to diverse technical and non-technical stakeholders. Generate patentable IP; publish and present research at leading venues (ICML, RLDM, CogSci, ICCC, DCC---Best Paper Award). 
      \end{cvitems}
    }

  \vspace{2mm}
  \cventry
    {Postdoctoral Scholar, Center for Cognitive Neuroscience}
    {Duke University}
    {Durham, NC}
    {Jun 2016 -- Jan 2021}
    {
      \begin{cvitems}
        \item Developed and validated multidisciplinary computational models of behavior change. Applied ML and hierarchical, spatial, time series, and multivariate statistical modeling to predict individual intervention response from large datasets containing survey, behavioral, eye-tracking, and brain imaging data. Published in top-tier journals including \textit{PNAS}, \textit{Nature Aging}, and \textit{Nature}.
        \item Secured competitive federal funding; managed research teams of 4--10 people. Mentored graduate students and research staff. 
      \end{cvitems}
    }

  \vspace{2mm}
  \cventry
    {Postdoctoral Associate, Institute of Cognitive Science}
    {University of Colorado, Boulder}
    {Boulder, CO}
    {Apr 2015 -- May 2016}
    {
      \begin{cvitems}
        \item Developed novel Bayesian approach for modeling variability in the brain’s processing of risky decisions. 
        \item Developed new incentive personalization model for shaping self-control in individual decision making.
      \end{cvitems}
    }

  \vspace{2mm}
  \cventry
    {Director, CaféWell Behavioral Labs}
    {Welltok, Inc.}
    {Denver, CO}
    {Aug 2013 -- Jan 2015}
    {
      \begin{cvitems}
        \item Owned the redesign of instrumentation for measuring user behavior, defining novel outcome metrics and rebuilding the data architecture for multivariate experimentation, translating requirements across engineering, product, and design into a unified, extensible measurement framework.
        \item Applied ML to user experience and optimized outcomes for individual users; translated algorithms into actionable designs using human-centered design principles.
        \item Led quantitative and qualitative UX research for mobile and web products; drove product design by synthesizing user testing, analytics, and market research.
      \end{cvitems}
    }

\end{cventries}

%-------------------------------------------------------------------------------
% EDUCATION
%-------------------------------------------------------------------------------
\vspace{2mm}
\cvsection{Education}

\begin{cventries}

  \cventry
    {Ph.D., Computation and Neural Systems}
    {Caltech}
    {Pasadena, CA}
    {2014}
    {}

  \cventry
    {B.A., Psychology}
    {Stanford University}
    {Stanford, CA}
    {2006}
    {}

\end{cventries}

%-------------------------------------------------------------------------------
% SKILLS
%-------------------------------------------------------------------------------
\vspace{2mm}
\cvsection{Skills}

\begin{cvskills}
  \cvskill
    {Technical}
    {Python, TensorFlow, R, Stan, MATLAB, bash, SQL}

  \cvskill
    {Methods}
    {Bayesian analysis and model comparison, forecasting and time series analysis, causal inference, reinforcement learning, state space models, sequential integrator models}
\end{cvskills}

%-------------------------------------------------------------------------------
\end{document}